\section{Background and Motivation}
\label{sec:background}

\subsection{Tool Routing}
\label{sec:tool-routing}

Tool routing is commonly formulated as a selection problem: given a user query
$q$ and a catalog of tools $\mathcal{T}$, choose a subset
$\{t_1,\ldots,t_k\}\subseteq\mathcal{T}$ to execute.

This formulation works for small catalogs that fit comfortably within an LLM's
context window. In production systems, however, tool ecosystems frequently
contain hundreds of available tools \cite{qin2023toolllm}. Presenting the full
catalog increases context length, degrades selection accuracy, raises inference
cost, and often produces hallucinated tool invocations
\cite{li2025toolscope}. When tasks require multiple tool calls, the model must
also infer a valid execution order while reasoning over the entire catalog.

These challenges suggest that tool routing is not only a retrieval problem, but
also a compositional reasoning problem.

\subsection{Existing Approaches}
\label{sec:existing-approaches}

\paragraph{Direct Selection.}

The simplest approach presents every tool description to the language model and
lets it decide which tools to invoke
\cite{schick2024toolformer,patil2023gorilla}. While effective for small tool
sets, this strategy scales poorly because context size grows linearly with the
number of tools and the model must simultaneously reason about tool semantics
and compatibility.

\paragraph{Retrieval-Based Routing.}

Embedding-based methods first retrieve a subset of semantically relevant tools
before invoking the language model \cite{li2025toolscope}. This substantially
improves scalability, but semantic similarity alone does not guarantee that the
retrieved tools can be composed into a valid workflow. Composition remains an
implicit responsibility of the LLM.

\paragraph{Graph-Based Planning.}

Several recent systems introduce graph structure to support tool composition.
PLaG \cite{lin2024plag} serializes task graphs into prompts for code
generation. ControlLLM \cite{zhang2024controlllm} constructs tool graphs and
searches for executable multimodal pipelines. GNN4Plan
\cite{wang2024gnn4plan} learns tool planning using graph neural networks, while
GRAFT \cite{li2025graft} incorporates graph representations through
graph-aware fine-tuning.

These approaches demonstrate that graph structure improves compositional
reasoning, but most require task-specific training, graph serialization inside
the prompt, or specialized planning architectures.

\subsection{Why Entities}
\label{sec:why-entities}

We argue that tool routing is formulated at the wrong level of abstraction.

Users naturally express tasks in terms of semantic entities rather than
executable tools. For example, the query

\begin{quote}
``What node is running my database pod?''
\end{quote}

mentions the entity types \texttt{Pod} and \texttt{Node} without referring to
any specific tool.

This observation suggests separating semantic reasoning from compositional
reasoning. Language models excel at mapping natural language to semantic
concepts, whereas graphs naturally encode structural constraints over valid
compositions. Rather than reasoning directly about tools, the language model
predicts entity types, and graph search recovers a valid sequence of tool
applications connecting them.

Formally, a tool registry induces a typed directed graph whose vertices
represent entity types and whose edges represent executable transformations.
Tool routing therefore becomes the problem of recovering a valid path between
predicted source and target entity types.

The benefit of this reformulation is not simply that there may be fewer entity
types than tools. In some domains, the two are nearly equal
(Section~\ref{sec:experimental-setup}). Instead, the graph itself constrains
the reasoning process. Once a target entity type is predicted, reverse
reachability immediately eliminates source types that cannot participate in any
valid composition. We quantify this structural reduction through entity pruning
(Section~\ref{sec:graph-constraint}), observing 56--87\% average pruning across
five domains, including domains where entity types and tools are nearly in
one-to-one correspondence.

This decomposition assigns each component the reasoning task for which it is
best suited. The language model performs semantic inference over entity types,
while graph search performs compositional reasoning over executable
transformations. As a result, valid tool sequences are recovered through graph
structure rather than inferred directly by the language model.